% Chapitre 2 - Page 10 | Travaux Avant-Projet | L3 ELN
\documentclass[11pt,a4paper]{article}
\usepackage{fontspec}
\setmainfont{TeXGyrePagella}
\usepackage[french]{babel}
\usepackage{graphicx,tikz,xcolor,geometry}
\usetikzlibrary{calc}
\geometry{margin=0pt}
\pagestyle{empty}
\definecolor{bleufonce}{HTML}{1B3A6B}
\definecolor{bleugris}{HTML}{5C7190}
\newcommand{\pos}[4]{\node[anchor=north west,inner sep=0pt,outer sep=0pt,text width=#3,align=left] at ($(current page.north west)+(#1,-#2)$) {#4};}
\newcommand{\sect}[1]{{\color{bleufonce}\bfseries\fontsize{17}{19}\selectfont #1}}
\newcommand{\step}[1]{{\color{bleufonce}\bfseries\fontsize{12.5}{14.5}\selectfont #1}}
\newcommand{\smallfont}{\fontsize{9.6}{11.8}\selectfont}
\newcommand{\legende}[1]{{\color{bleugris}\fontsize{8.5}{10}\selectfont #1}}
\newcommand{\puce}{\textcolor{bleufonce}{\textbullet}\enspace}
% \shot{x}{y}{largeur}{hauteur}{fichier} : capture encadrée
\newcommand{\shot}[5]{\node[anchor=north west,inner sep=0pt,outer sep=0pt] at ($(current page.north west)+(#1,-#2)$) {\includegraphics[width=#3]{#5}};
  \draw[bleufonce!75,line width=.55pt,rounded corners=1mm]
    ($(current page.north west)+(#1,-#2)+(-.4mm,.4mm)$) rectangle ($(current page.north west)+(#1,-#2)+(#3,-#4)+(.4mm,-.4mm)$);}
\begin{document}
\begin{tikzpicture}[remember picture,overlay]
\draw[bleufonce,line width=.95pt,rounded corners=8mm]
 ($(current page.north west)+(5mm,-5mm)$) rectangle ($(current page.south east)+(-5mm,5mm)$);

\pos{10mm}{13mm}{189mm}{\sect{7. Les vues Schematic et PCB, puis l'exportation}}
\pos{10mm}{25mm}{188mm}{{\fontsize{10.5}{13}\selectfont Une fois le montage réalisé sur la plaque d'essai, il suffit de changer d'onglet : Fritzing affiche \textbf{automatiquement} le même circuit sous forme de schéma puis de circuit imprimé.}}

% --- Étape 5 : Schematic -------------------------------------------------------
\pos{10mm}{39mm}{188mm}{\step{Étape 5 --- La vue Schematic (schéma électrique)}}
\shot{14mm}{48mm}{112mm}{58.2mm}{assets/F06_schematic.png}
\pos{133mm}{49mm}{64mm}{{\smallfont
  \puce Les composants sont remplacés par leurs \textbf{symboles}.\\[1.4mm]
  \puce Les \textbf{lignes fines} (\emph{ratsnest}) indiquent les liaisons venant de la vue Breadboard.\\[1.4mm]
  \puce Le bouton \textbf{Autoroute} trace les fils automatiquement ; on peut ensuite les déplacer pour rendre le schéma lisible.}}
\pos{14mm}{108mm}{184mm}{\legende{Figure 11 --- La vue \textbf{Schematic} : l'Arduino Uno, la photorésistance sur A0 et la LED sur la broche 2.}}

% --- Étape 6 : PCB -----------------------------------------------------------------
\pos{10mm}{118mm}{188mm}{\step{Étape 6 --- La vue PCB (circuit imprimé)}}
\shot{14mm}{127mm}{112mm}{58.9mm}{assets/F07_pcb.png}
\pos{133mm}{128mm}{64mm}{{\smallfont
  \puce La carte a ici la forme d'un \textbf{shield} Arduino.\\[1.4mm]
  \puce Les \textbf{pistes} relient les pastilles des composants ; \textbf{Both Layers} permet de travailler sur les deux faces.\\[1.4mm]
  \puce \textbf{Autoroute} trace les pistes ; le message \textbf{Routing completed} confirme que toutes les liaisons sont faites.}}
\pos{14mm}{188mm}{184mm}{\legende{Figure 12 --- La vue \textbf{PCB} : placement des composants et tracé des pistes.}}

% --- Étape 7 : Export ------------------------------------------------------------
\pos{10mm}{198mm}{188mm}{\step{Étape 7 --- Exporter le travail}}
\shot{14mm}{207mm}{62mm}{47.3mm}{assets/F08_export.png}
\pos{83mm}{208mm}{114mm}{{\smallfont Le menu \textbf{File \(\rightarrow\) Export} propose :\\[1.6mm]
  \puce \textbf{as Image} : image (PNG, SVG) ou PDF de la vue active, pour un rapport ;\\[1.2mm]
  \puce \textbf{for Production} : fichiers \textbf{Gerber} pour faire fabriquer le PCB ;\\[1.2mm]
  \puce \textbf{List of parts (Bill of Materials)} : la liste des composants à acheter ;\\[1.2mm]
  \puce \textbf{XML / SPICE Netlist} : la liste des connexions, pour d'autres logiciels.}}
\pos{14mm}{256.5mm}{184mm}{\legende{Figure 13 --- Le menu \textbf{File \(\rightarrow\) Export}.}}

\draw[bleufonce,line width=.7pt,rounded corners=2mm]
 ($(current page.north west)+(10mm,-263mm)$) rectangle ($(current page.north west)+(200mm,-276mm)$);
\pos{15mm}{265.2mm}{180mm}{{\fontsize{9.3}{11.2}\selectfont \textcolor{bleufonce}{\bfseries À retenir :} on dessine une seule fois dans Breadboard ; Schematic et PCB se complètent, puis \textbf{Export} fournit les documents du projet.}}

\draw[bleufonce!55,line width=.45pt] ($(current page.south west)+(10mm,14.7mm)$)--($(current page.south east)+(-10mm,14.7mm)$);
\node[anchor=south west,inner sep=0pt,text=bleufonce!80] at ($(current page.south west)+(10mm,7.4mm)$) {\fontsize{9.5}{11}\selectfont 2026/2027};
\node[anchor=south east,inner sep=0pt,text=bleufonce!80] at ($(current page.south east)+(-10mm,7.4mm)$) {\fontsize{9.5}{11}\selectfont Page 10};
\end{tikzpicture}
\null
\end{document}
